\Needspace{22\baselineskip}
\section*{第17章：子类型化的 ML 实现}
\addcontentsline{toc}{section}{第17章：子类型化的 ML 实现}

\begin{taplsolutionentrybox}
\phantomsection
\label{sol:translator:17.3.3}
\textbf{练习17.3.3}

下面的实现让子类型检查失败时携带两类信息：比较进行到记录的哪一层，以及
该处究竟是缺少字段，还是两侧类型的构造不同。\texttt{subtype\_error}
保存这些数据，\texttt{Subtype\_error} 则把数据包装成可由 \texttt{raise}
抛出的异常。这样既保留了结构化错误信息，也符合 OCaml 对异常值类型的要求。

\texttt{typeof} 完整处理记录、投影、变量、抽象和应用。应用分支先分别求出
函数与实参的类型，再从两个类型的最外层开始调用 \texttt{check\_subtype}；
递归检查进入记录字段或函数类型内部时，会逐步记录所经过的位置。投影和变量
分支也报告缺失字段、越界索引及上下文长度。错误消息会打印出错路径，以及实际
类型和预期类型。完整源码还实现错误消息格式化，并用测试覆盖成功检查、嵌套字段
缺失、类型构造不匹配、变量越界、投影缺少字段、投影非记录项和应用非函数项。
下面省略本章正文已经展示过的项、类型和上下文定义，以及结构未变的类型检查
分支，只列出错误数据、带路径的子类型检查和应用分支。

\taplocamlfileflowchunkfirst{../code/ocaml/chapter17/diagnostics.ml}{13}{26}
\taplocamlfileflowchunk{../code/ocaml/chapter17/diagnostics.ml}{28}{48}
\taplocamlfileflowchunk{../code/ocaml/chapter17/diagnostics.ml}{69}{79}

Rust 对照实现沿用本章正文的 \texttt{Type}、\texttt{Term} 和
\texttt{Context}。三个代码段分别给出诊断数据、带路径的子类型检查，以及把
诊断接入函数应用类型检查的入口。

\taplrustsupportflowchunkfirst{sol-translator-17-diagnostic-error}{1}{30}
\taplrustsupportflowchunk{sol-translator-17-diagnostic-subtype}{1}{20}
\taplrustsupportflowchunk{sol-translator-17-diagnostic-subtype}{21}{36}
\taplrustsupportflowchunk{sol-translator-17-diagnostic-app}{1}{20}

\taplsolutionbacklink{ex:17.3.3}{返回练习17.3.3}
\end{taplsolutionentrybox}

\begin{taplsolutionentrybox}
\phantomsection
\label{sol:translator:17.3.4}
\textbf{练习17.3.4}

强制转换语义把隐含的子类型使用编译成普通函数调用。先区分源语言类型与
目标语言类型。类型翻译记作 \(\llbracket T\rrbracket\)，其中

\[
  \llbracket\tapltype{Top}\rrbracket = \tapltype{Unit},\qquad
  \llbracket S\to T\rrbracket
    = \llbracket S\rrbracket\to\llbracket T\rrbracket,\qquad
  \llbracket\{l_i:T_i\}^{i\in1..n}\rrbracket
    = \{l_i:\llbracket T_i\rrbracket\}^{i\in1..n}
\]

因此，\texttt{coerce S T} 在确认 \(S\subtype T\) 后，生成的函数具有类型
\(\llbracket S\rrbracket\to\llbracket T\rrbracket\)，并非未经翻译的
\(S\to T\)。特别地，目标为 \(\tapltype{Top}\) 时生成
\(\lambda x:\llbracket S\rrbracket.\,\texttt{unit}\)，其结果确实具有
\(\tapltype{Unit}\) 类型。函数类型的参数位置反向生成强制转换，结果位置沿原方向
生成；记录类型则只重建目标类型要求的字段。

例如，把
\[
  \{x=\{\},y=\{\}\}
\]
传给只接受 \(\{x:\tapltype{Top}\}\) 的函数时，类型检查使用记录宽度
子类型化。翻译器会插入一个显式转换函数，其目标语言形式为
\[
  \lambda r:\{x:\tapltype{Unit},y:\tapltype{Unit}\}.\,\{x=r.x\}
\]
这个函数丢弃字段 \(y\)，把实参变成函数实际要求的记录。翻译后的程序因而
不再需要子类型规则。

项翻译同时返回目标类型与目标项。应用 \(t_1\ t_2\) 若得到源类型
\(S_1\to S_2\) 和 \(T_1\)，便在实参与函数之间插入从
\(\llbracket T_1\rrbracket\) 到 \(\llbracket S_1\rrbracket\) 的强制转换。
抽象中的类型标注、上下文中保存的类型以及记录字段类型也都经过同一类型翻译。

下面只列出本题新增的核心 OCaml 实现：源语言与目标语言的语法、类型翻译、
强制转换构造及项翻译。目标语言的移位、替换与求值沿用\chapref{7}的算法，
完整配套源码中保留这些实现和端到端测试。

\tcbbreak\newpage
\taplocamlfileflowchunkfirst{../code/ocaml/chapter17/coercion.ml}{1}{28}
\taplocamlfileflowchunk{../code/ocaml/chapter17/coercion.ml}{29}{56}
\taplocamlfileflowchunk{../code/ocaml/chapter17/coercion.ml}{57}{73}
\taplocamlfileflowchunk{../code/ocaml/chapter17/coercion.ml}{100}{120}
\taplocamlfileflowchunk{../code/ocaml/chapter17/coercion.ml}{121}{139}

Rust 对照实现采用相同结构：

\tcbbreak\newpage
\taplrustsupportflowchunkfirst{sol-translator-17-coercion-support}{1}{25}
\taplrustsupportflowchunk{sol-translator-17-coercion-support}{26}{47}
\taplrustsupportflowchunk{sol-translator-17-coercion}{1}{19}
\taplrustsupportflowchunk{sol-translator-17-coercion}{20}{38}
\taplrustsupportflowchunk{sol-translator-17-coercion}{39}{55}
\taplrustsupportflowchunk{sol-translator-17-coercion}{56}{78}
\taplrustsupportflowchunk{sol-translator-17-coercion}{79}{103}
\taplrustsupportflowchunk{sol-translator-17-translate}{1}{18}
\taplrustsupportflowchunk{sol-translator-17-translate}{19}{33}

端到端测试构造上述宽记录应用，检查翻译后的项仍为良类型，并最终求值为
\texttt{unit}：

\taplrustsupportflowchunk{sol-translator-17-example}{1}{20}

\taplsolutionbacklink{ex:17.3.4}{返回练习17.3.4}
\end{taplsolutionentrybox}
